\documentclass[dvisvgm,tikz,border=3pt]{standalone}
\usepackage{amsmath,amssymb}
\usepackage{pgfplots}
\usepackage{CJKutf8}
\pgfplotsset{compat=1.18}
\usetikzlibrary{arrows.meta,calc,positioning}

\definecolor{themebg}{HTML}{2e362c}
\definecolor{themefg}{HTML}{edf4e8}
\definecolor{themegray}{HTML}{a4af9d}
\definecolor{themecurve}{HTML}{7eb6ff}
\definecolor{themealert}{HTML}{ff7b7b}
\definecolor{themeamber}{HTML}{f5b942}

\pgfdeclarelayer{background}
\pgfdeclarelayer{foreground}
\pgfsetlayers{background,main,foreground}

\begin{document}
\begin{CJK*}{UTF8}{gbsn}
\pagecolor{themebg}
\begin{tikzpicture}[color=themefg, text=themefg]

\begin{axis}[
    width=13.0cm,
    height=8.5cm,
    axis lines=middle,
    axis line style={color=themefg, -{Stealth}},
    tick style={color=themefg},
    tick label style={color=themefg, font=\scriptsize},
    every axis label/.style={color=themefg},
    xmin=-2.8, xmax=3.2,
    ymin=-3.5, ymax=3.5,
    xtick={-1, 1},
    xticklabels={$-1$, $1$},
    ytick={-2, 2},
    yticklabels={$-2$, $2$},
    clip=false,
    xlabel={$x$},
    ylabel={$y$},
    title style={font=\footnotesize\bfseries, at={(0.5,1.06)}, anchor=south},
    title={方程根数与水平扫描线相交：极值点切触为根数突变的临界边界},
]

% 1. 三次多项式主曲线 y = x^3 - 3x
\addplot[themecurve, line width=2.2pt, domain=-2.3:2.3, samples=80] {x^3 - 3*x};
\node[text=themecurve, font=\scriptsize\bfseries, anchor=west] at (axis cs:1.8, 2.8) {$y=x^3-3x$};

% 2. 极大值与极小值水平切线 (临界高度 k = 2 与 k = -2)
\draw[themealert, dashed, line width=1.3pt] (axis cs:-2.6, 2) -- (axis cs:2.8, 2);
\node[text=themealert, font=\scriptsize\bfseries, anchor=west] at (axis cs:2.8, 2) {$k=2$ (2 根: 1切1穿)};

\draw[themealert, dashed, line width=1.3pt] (axis cs:-2.6, -2) -- (axis cs:2.8, -2);
\node[text=themealert, font=\scriptsize\bfseries, anchor=west] at (axis cs:2.8, -2) {$k=-2$ (2 根: 1切1穿)};

% 3. 3 根区间水平扫描线 (如 k = 0.5)
\draw[themeamber, line width=1.4pt] (axis cs:-2.6, 0.5) -- (axis cs:2.8, 0.5);
\node[text=themeamber, font=\scriptsize\bfseries, anchor=west] at (axis cs:2.8, 0.5) {$k \in (-2,2)$ (3 个不同实根)};

% 4. 1 根区间水平扫描线 (如 k = 2.8)
\draw[themegray, line width=1.2pt] (axis cs:-2.6, 2.8) -- (axis cs:2.8, 2.8);
\node[text=themegray, font=\scriptsize\bfseries, anchor=west] at (axis cs:2.8, 2.8) {$k > 2$ (1 个实根)};

% 5. 极值点与相交实根标记
\begin{pgfonlayer}{foreground}
    % 极大值点 (-1, 2)
    \addplot[only marks, mark=*, mark size=3.0pt, fill=themealert, draw=themefg] coordinates {(-1, 2)};
    \node[text=themealert, font=\scriptsize\bfseries, anchor=south east] at (axis cs:-1.05, 2.05) {极大值 $(-1, 2)$};

    % 极小值点 (1, -2)
    \addplot[only marks, mark=*, mark size=3.0pt, fill=themealert, draw=themefg] coordinates {(1, -2)};
    \node[text=themealert, font=\scriptsize\bfseries, anchor=north west] at (axis cs:1.05, -2.05) {极小值 $(1, -2)$};

    % k=0.5 处的 3 个交点 (暖金圆点)
    \addplot[only marks, mark=*, mark size=2.5pt, themeamber] coordinates {(-1.64, 0.5) (-0.17, 0.5) (1.81, 0.5)};

    % 结论框放置在左上空白区
    \node[text=themefg, font=\scriptsize, anchor=north west] at (axis cs:-2.7, 3.4) {
        \begin{tabular}{l}
        \textbf{根数分类结论}：\\
        1. $k > 2$ 或 $k < -2$：恰有 \textbf{1 个实根}\\
        2. $k = 2$ 或 $k = -2$：恰有 \textbf{2 个实根} (含二重切根)\\
        3. $-2 < k < 2$：恰有 \textbf{3 个不同实根}
        \end{tabular}
    };
\end{pgfonlayer}

\end{axis}
\end{tikzpicture}
\end{CJK*}
\end{document}
